\begin{center}
\thispagestyle{fancy}
\LARGE{\textbf{Abstract}}
\end{center}
\linespread{1.33}
\large{
\paragraph{}
Property tax collection in Indian municipalities remains largely manual, inefficient, and prone to revenue leakage. Most smaller Urban Local Bodies continue to rely on paper registers, inconsistent assessment practices, and outdated property records. Citizens face difficulties understanding tax calculations and making timely payments, while municipal officials lack the tools needed to track compliance and identify defaulters effectively.

\paragraph{}
The Smart Property Tax Management System (SMPTMS) is a full-stack web application developed specifically to address these challenges for Jaipur Nagar Nigam. The platform digitizes the complete property tax lifecycle — from property registration and automated tax assessment to online payment processing and official receipt generation. Citizens can register their properties, view computed tax amounts based on plot area, zone, and land use type, and complete payments through an integrated Razorpay gateway, all without visiting a municipal office.

\paragraph{}
On the administrative side, officials access a dedicated dashboard featuring a GIS-based property map powered by Leaflet.js, showing all registered properties across Jaipur with real-time payment status indicators. Zone-wise analytics, defaulter tracking, and exportable reports provide actionable insights for revenue management and urban planning.

\paragraph{}
The system was developed using Next.js 14 (App Router), TypeScript, Supabase (PostgreSQL), Tailwind CSS, and Razorpay, and is deployed live at smptms.vercel.app. The tax calculation follows a standardized formula: Tax = (Plot Area x DLC Rate) / 2000, eliminating manual manipulation and ensuring transparency. SMPTMS demonstrates that modern web technologies can deliver a lightweight, cost-effective, and genuinely usable municipal tax platform for cities like Jaipur.
}
